\begin{theorem}[E.32 --- Proved: full pronilpotent sections at the three- and five-power levels]
\label{E-full-pronilpotent}
Let $n=\ell^r$ with $\ell\in\{3,5\}$ and $r\ge1$, and let
$N=\pi_1(U_{n,\overline{\mathbb Q}})$ be the full profinite
geometric group of the plus Fermat curve. Use the rational
unit tangent at $O=[0:-1:1]$ to specify the arithmetic action.
Its maximal pronilpotent quotient has the canonical description
\[
 N^{\mathrm{nil}}=\prod_q N^{(q)},
 \qquad
 N/\gamma_{c+1}N\simeq
       \prod_q N^{(q)}/\gamma_{c+1}N^{(q)}\quad(c\ge1).
 \tag{E.32.1}
\]
Here $q$ ranges over all primes and lower central subgroups
are closed. At each quotient, the cuspidal classes in this
statement are the images of the actual cusp packets under
$D_c\hookrightarrow\pi_1(U_n)\twoheadrightarrow
E^{(c)}$. This specifies the packet map even at class one,
where the normalizer of a boundary line can contain further
sections. The arithmetic extension with geometric kernel
$N^{\mathrm{nil}}$ has infinitely many continuous section
classes whose image in every displayed positive-class
nilpotent quotient is noncuspidal. Consequently no level of
this full profinite nilpotent tower selects precisely the
cuspidal section classes at these Fermat levels.
\end{theorem}
\begin{proof}
Every finite nilpotent group is the direct product of its
Sylow subgroups. A map from $N$ to such a group factors,
in each Sylow coordinate, through $N^{(q)}$. Conversely,
the natural map from $N$ to a finite product of its finite
prime-power quotients is onto: its image is a subdirect
product, and factors of coprime orders have no common
nontrivial quotient. Compactness now proves that
$N\to\prod_qN^{(q)}$ is onto. This proves the first
identity in (E.32.1), with its actual quotient map.

The quotient $N/\gamma_{c+1}N$ is the inverse limit of
its finite nilpotent quotients of class at most $c$.
Apply the preceding factorization to each such quotient.
The reverse comparison is induced by the coordinate
maps to $N^{(q)}/\gamma_{c+1}N^{(q)}$, whose product has
class at most $c$. The universal properties give inverse
continuous maps, proving the second identity. Both
comparisons commute with every automorphism of $N$,
and in particular with the specified Galois action.

Write a section supplied by E.16 as
$s_\ell(g)=b_\ell(g)s_O(g)$, where
$b_\ell:G_{\mathbb Q}\to N^{(\ell)}$ is its actual
continuous nonabelian cocycle. Define a cocycle in the
first product of (E.32.1) by taking coordinate
$b_\ell$ at $\ell$ and the identity at every other prime.
The product topology and the coordinatewise arithmetic
action make it a continuous cocycle. It therefore gives
one actual section of the pronilpotent arithmetic
extension, whose images are compatible in every
finite nilpotent quotient involving any finite set
of coefficient primes.

Choose the E.16 sections with distinct infinite-order
images in
$H^1(\mathbb Q,T_\ell\operatorname{Jac}(X_n))$.
That compactified abelian image factors through every
quotient in (E.32.1). Every cusp image is $n$-torsion:
the ratio of two cusp tangent-line equations has
divisor $n(d-O)$, as checked in E.16. Thus none of the
chosen images can belong to a cusp packet at any
positive class. Distinct compact images also prevent
conjugacy of the corresponding sections. This proves
the assertions, including the full coefficient-prime
comparison. It does not compute section classes of
finite quotients outside the specified nilpotent
tower. The anchors are \S51.16, folios733v--736v,
scans172--178, and LM7-E.23(iii).
\end{proof}

\begin{theorem}[E.33 --- Proved: full metabelian zero-abelian families at every three-power level]
\label{E-full-metabelian-tower}
Let $n=3^r$ with $r\ge1$ and keep the full profinite
group $N$ and tangent $s_O$ of E.32. Set
\[
 M=N/N'',\qquad A=N^{\mathrm{ab}},\qquad
 d_n=\frac{(n^2-1)(n^2+3)}4.
\]
There is a continuous injective family
\[
 \mathbb Z_3^{d_n}\longrightarrow
 \operatorname{Sec}(M\rtimes_{s_O}G_{\mathbb Q})/M
 \tag{E.33.1}
\]
all of whose abelianized classes equal $[s_O]$.
The zero parameter maps to $[s_O]$; every nonzero
parameter is noncuspidal. In particular $d_3=24$.
The family is specified by a choice of a free lattice
in the class-two cohomology quotient below and lifts
of its generators in the pro-three derived module.
E.31 supplies a separate explicit finite certificate
for one cubic direction.
\end{theorem}
\begin{proof}
The geometric rank is $n^2+1$. Put
\[
 P=N^{(3)},\quad V=P^{\mathrm{ab}},\quad
 D_3=P'/P'',\quad B_2=\bigwedge^2_{\mathbb Z_3}V.
\]
The continuous surjection $D_3\twoheadrightarrow B_2$
is the degree-two commutator map, with kernel
$\gamma_3P/P''$. The Fox and Haar comparison of E.31
works in every finite free rank, by the identical
row $\sum r_i(X_i-1)$. It gives the actual
arithmetic-equivariant embedding
\[
 j:D_3\hookrightarrow N'/N'',
 \qquad qj=1,
 \tag{E.33.2}
\]
where $q$ is the natural pro-three comparison.
No change in the arithmetic action is made.

By E.16 that action on $P$ factors through
\[
 \Gamma=\langle\tau,\tau'\rangle_{\mathrm{pro}-3}
                 \rtimes\langle c\rangle,
 \qquad c^2=1,\quad c\tau c=\tau'.
 \tag{E.33.3}
\]
Use the conjugate Frobenius basis at seven; let
$T$ be the action of $\tau$ on $B_2$. Since two is
invertible, every $B_2$-cocycle class has a
representative with $b(c)=0$. Such a representative
is specified by its arbitrary value at $\tau$;
the value at $\tau'$ must be $c b(\tau)$. Freeness
extends these assignments continuously, and they
respect the two relations involving $c$. A
coboundary preserves $b(c)=0$ precisely when its
parameter belongs to $B_2^c$. Consequently
\[
 H^1_{\mathrm{cts}}(\Gamma,B_2)
       =B_2/(T-1)B_2^c.
 \tag{E.33.4}
\]
This is a quotient by the literal compact image.

The eigenvalues of good Frobenius on the compact
part of $V$ have weight one and those on the
boundary part have weight two. All eigenvalues
on $B_2$ therefore have weights two, three or
four, so none equals one. Thus $T-1$ is invertible
on $B_2\otimes\mathbb Q_3$, and (E.33.4) has
free rank $\dim(B_2\otimes\mathbb Q_3)^-$.

We compute that dimension with the actual real
action. The compact curve has genus
$g=(n-1)(n-2)/2$; complex conjugation has
$g$ eigenvalues of each sign on its first
homology. This follows from its reversal of
the nondegenerate intersection pairing: the
two eigenspaces are isotropic, and the pairing
identifies either with the dual of the other.
Of the $3n$ boundary points, exactly three
are real, and the rest form $(3n-3)/2$
conjugate pairs. The boundary module is the
quotient of their permutation module tensored
with $\mathbb Z_3(1)$ by the diagonal line
generated by the sum of all cusp generators.
Complex conjugation acts by minus one on that
line; thus taking this quotient removes one
minus eigenvalue.
Its plus and minus dimensions are respectively
$(3n-3)/2$ and $(3n+1)/2$. Hence
\[
 \dim V^+=\frac{n^2-1}2,
 \qquad
 \dim V^-=\frac{n^2+3}2,
 \qquad
 \dim B_2^-=\dim V^+\dim V^-=d_n.
 \tag{E.33.5}
\]
The last equality uses the actual direct sum
$\bigwedge^2(V^+\oplus V^-)$: its minus part is
$V^+\otimes V^-$. Thus no trace estimate replaces
the eigenspace calculation.
The program \texttt{metabelian\_family\_rank.py}
constructs the cubic boundary quotient matrix
and its full exterior-square action independently;
its exact nullspace ranks are $4$, $6$ and $24$.

Choose $d_n$ elements of $B_2$ whose classes
give a free lattice in (E.33.4), and choose
their lifts $d_i\in D_3$. For
$a=(a_1,\ldots,a_{d_n})\in\mathbb Z_3^{d_n}$
assign
\[
 b_a(\tau)=\sum_i a_i d_i,\qquad
 b_a(\tau')=c b_a(\tau),\qquad b_a(c)=0.
 \tag{E.33.6}
\]
The same free-source argument gives an actual
continuous $D_3$-cocycle. Its dependence on $a$
is continuous: equality on the free generators
gives $b_a=\sum_i a_i b_i$ for the fixed cocycles
$b_i$. This is jointly continuous in $(a,g)$,
and therefore continuous in the compact-open
topology. Inflate to $G_{\mathbb Q}$
and set $s_a(g)=j(b_a(g))s_O(g)$. This proves
existence of (E.33.1) with zero abelian cocycle.

The invariant group $H^0(G_{\mathbb Q},A)$ is
zero, coefficient prime by coefficient prime,
using a good Frobenius of different residue
characteristic and the same compact and
boundary weights. Thus a geometric conjugator
between two such sections lies in $N'/N''$.
Projecting with $q$ and then to $B_2$ proves
that their difference is a coboundary in
(E.33.4). Inflation from $\Gamma$ creates no
new coboundaries: the actual action factors
through $\Gamma$, so the same parameter
defines a coboundary before and after
inflation. The chosen free lattice therefore
proves injectivity. Compactness and the
closed conjugacy relation give the stated
topology on this family of section classes.

We check cusp normalizers and packets in the
quotient itself. A cusp generator can be
included in a free basis of $N$, by omitting
a different cusp in the surface presentation.
For such a generator $x$, class two shows
that the abelian coordinate of any element
centralizing $x$ is a multiple of $x$.
After subtracting that power, the element
lies in $D=N'/N''$ and is killed by $X-1$.
The Fox embedding $D\hookrightarrow R^{n^2+1}$
and injectivity of $X-1$ on $R$ force it
to vanish. For the last injectivity, in each
coefficient-$q$ factor refine a finite
abelian quotient by a further $q^m$ in the
primitive $X$ direction. An $X$-invariant
coefficient pushes forward as a multiple
of $q^m$, and is zero modulo every fixed
$q$-power once $m$ is large. The inverse
limit is therefore zero. Consequently
$C_M(x)=\langle x\rangle$. A normalizer
cannot change the exponent of $x$, since
conjugation fixes its primitive abelian
vector. Thus its geometric normalizer is
the same cusp group. The arithmetic
normalizer over a rational cusp has
exactly the usual Kummer packet. Over a
nonrational cusp its image is the
stabilizer of that cusp: distinct boundary
lines remain distinct in $A$ and the
Galois action permutes them. Such a
normalizer cannot contain a $G_{\mathbb Q}$
section.

For the three rational cusps, the
tangent-line ratio has divisor $n(d-O)$.
The Abel--Jacobi image of $d-O$ is nonzero
for $d\ne O$, by injectivity of the curve's
Abel map in genus at least one. It is
therefore nonzero three-primary torsion.
Its three-primary Kummer class is nonzero
by the Mordell--Weil Kummer injection.
No packet at $d\ne O$ has zero abelian
class. At $O$ the primitive cusp line $I_O$
and $H^0(G_{\mathbb Q},A/I_O)=0$ make the
packet map injective on abelianization.
Its zero class is precisely $s_O$.
Injectivity of (E.33.1) excludes this
class for every nonzero parameter.
The anchors are \S51.16, folios733v--736v,
scans172--178, and LM7-E.23(iii).
\end{proof}
